% Generated draft: requires mathematical and source review.
\documentclass[11pt]{article}
\usepackage[margin=25mm]{geometry}
\usepackage{fontspec}
\setmainfont{Latin Modern Roman}
\usepackage{amsmath,amssymb,mathtools,unicode-math}
\setmathfont{Latin Modern Math}
\usepackage{enumitem,needspace,xurl,hyperref}
\setlength{\parindent}{0pt}
\setlength{\parskip}{5pt}
\setlength{\emergencystretch}{3em}
\newcommand{\N}{\mathbb{N}}
\newcommand{\Z}{\mathbb{Z}}
\newcommand{\Q}{\mathbb{Q}}
\newcommand{\R}{\mathbb{R}}
\newcommand{\C}{\mathbb{C}}
\newcommand{\F}{\mathbb{F}}
\newcommand{\A}{\mathbb{A}}
\newcommand{\PP}{\mathbb{P}}
\newcommand{\E}{\mathbb{E}}
\newcommand{\eps}{\varepsilon}
\newcommand{\dd}{\,\mathrm{d}}
\begin{document}
\section*{UnsolvedMath ID: 800948}
\textbf{Source problem code:} catalogue entry
\subsection{Definitions and mathematical statement}
The supplied extract contains a closely related question, but omits the preceding question and its hypotheses. To express the displayed limits, let \(d\) denote the intended ambient distance and let \(\alpha(t)\) and \(\beta(t)\) denote the intended curves, with parameters extending to both \(-\infty\) and \(+\infty\). The question is whether
\[
\lim_{t\to-\infty}d\bigl(\alpha(t),\beta(t)\bigr)=0
\quad\Longrightarrow\quad
\lim_{t\to+\infty}d\bigl(\alpha(t),\beta(t)\bigr)=+\infty.
\]
Here the conclusion means that for every \(M>0\) there exists \(T\in\R\) such that \(d(\alpha(t),\beta(t))>M\) for every \(t>T\) in the common parameter domain. The intended ambient space, metric hypotheses, class of geodesics, parametrization, and any required distinctness of \(\alpha\) and \(\beta\) are not supplied. Thus this is an incomplete question template, not an assertion quantified over arbitrary metric spaces and curves.

The surrounding source text reports that the answer to the preceding question is positive in dimension \(2\), citing [Green1954], and open in higher dimensions; it also reports that the proof in [Green1956] is incorrect. Because the preceding question is absent, these qualifications cannot reliably be assigned to the displayed implication. The same text reports a negative answer to the corresponding question for Finsler metrics: the Teichm\"uller metric on Teichm\"uller space has no conjugate points, while Masur exhibited two rays through the same point which do not diverge [Masur1975]. The extract does not specify the precise meaning of divergence in that preceding question.
\subsection{Short English statement}
The extracted question asks whether two intended geodesics whose distance tends to zero in backward time must have distance tending to infinity in forward time. Essential hypotheses are missing, so a complete quantified formulation cannot be recovered. The supplied registry labels the record open; its historical and literature claims require review.
\subsection{Sources}
\begin{itemize}
\item Local UnsolvedMath registry, permanent ID \(800948\), source code \texttt{catalogue entry}, title \emph{Open Problems and Questions about Geodesics --- Question 28}. The record attributes the problem to Burns and Matveev, with proposed year 2019, and supplies status \texttt{open}.
\item Burns and Matveev, \emph{Open Problems and Questions about Geodesics} (2019, as identified in the supplied source list), Question 28: \url{https://arxiv.org/abs/1308.5417}. The supplied provenance records access on 2026-10-16 and extraction method \texttt{arxiv\_latex}; this access information is supplied provenance, not an independently verified visit.
\item [Green1954], cited in the supplied extract for the positive dimension-\(2\) answer to the omitted preceding question. No full bibliographic entry is supplied.
\item [Green1956], whose proof is described as incorrect in the supplied extract. No full bibliographic entry is supplied.
\item [Masur1975], cited in the supplied extract for two rays through the same point in Teichm\"uller space which do not diverge. No full bibliographic entry is supplied.
\end{itemize}
\end{document}
